% Zumdahl Chapter 17 -- Spontaneity, Entropy, and Free Energy.
% Elaborated notes, 5 blocks. Covers CED 9.1-9.7 (front half of Unit 9).
\documentclass{apchem-notes}

\apcunitword{Chapter}
\apcunit{17}
\apcunittitle{Entropy and Free Energy}
\apcdoctitle{Guided Notes}
\apcsubtitle{Zumdahl \S17.1--17.10 \textbullet\ PDF pp.~834--880 \textbullet\ 5 blocks}

\begin{document}
\apctitleblock

\begin{apcnote}[How this chapter fits the AP course]
This is the \textbf{front half of Unit 9}, the last unit of the course:
\S17.1--17.2 $\to$ CED 9.1, \S17.6 $\to$ 9.2, \S17.3--17.4 and 17.7
$\to$ 9.3, \S17.9 $\to$ 9.5, \S17.5 $\to$ 9.6, \S17.10 $\to$ 9.7.
Chapter 18 supplies the electrochemistry half (9.8--9.11).

Kinetic control (\textbf{9.4}) has no Zumdahl section of its own --- it is a
deliberate callback to Chapter 12, and Block 4 makes it.

\S17.8 (pressure dependence of $G$) and the derivations in \S17.10 go
beyond the framework; the parts you need are folded into Blocks 4 and 5.

\textbf{This chapter answers a question Chapter 12 could not.} Kinetics told
you \emph{how fast}; equilibrium told you \emph{how far}. Neither told you
\emph{why} a reaction runs one way at all. That is what free energy is for.
\end{apcnote}

\begin{apctrap}
\textbf{Read this before anything else: the textbook and the exam use
different words.}

Zumdahl's chapter is titled ``Spontaneity, Entropy, and Free Energy'' and
uses \term{spontaneous} on nearly every page. The 2024 CED explicitly
retires that word:

\begin{quote}\small
``Historically, the term \emph{spontaneous} has been used to describe
processes for which $\Delta G^\circ < 0$. The phrase \emph{thermodynamically
favored} is preferred instead, so that common misunderstandings (equating
`spontaneous' with `suddenly' or `without cause') can be avoided.''
\end{quote}

Read ``spontaneous'' in the textbook; write \textbf{``thermodynamically
favored''} on the exam. They mean the same thing --- $\Delta G^\circ < 0$
--- and neither means fast, sudden, or uncaused. A thermodynamically favored
reaction may take a million years.
\end{apctrap}

% =====================================================================
\blocklesson{Entropy and the Second Law}{Zumdahl \S17.1--17.2}

\begin{objectives}
  \item Identify the sign and relative magnitude of an entropy change.
  \item State the second law and explain what it predicts.
\end{objectives}

\reading{Zumdahl \S17.1--17.2}{835--843}

\begin{warmup}[4]
  \item Sign of $\Delta H$ for an exothermic reaction:
        \answerblank[2cm]{negative}
  \item Which phase has the most freedom of motion?
        \answerblank[2cm]{gas}
  \item At constant $P$, $q$ equals: \answerblank[2.4cm]{$\Delta H$}
  \item Raising the temperature of a gas does what to the spread of
        molecular speeds? \answerblank[3cm]{broadens it}
\end{warmup}

\segment{INSTRUCTION A}{25}{What entropy actually measures}

\notesectionz{Dispersal, not disorder}{17.1}
\practice{6}

\term{Entropy} ($S$) measures how \textbf{dispersed} the matter and energy
of a system are --- how many microscopic arrangements produce the same
macroscopic state. The more ways there are, the higher the entropy.

``Disorder'' is the usual shorthand and it is a poor one: a messy bedroom is
not a thermodynamic state. Think \emph{spread out}, in two senses the CED
names separately:

\begin{itemize}[leftmargin=*,itemsep=0.3em]
  \item \textbf{Matter disperses.} Particles occupy a larger volume or move
        more freely.
  \item \textbf{Energy disperses.} At higher temperature the distribution of
        molecular kinetic energies \answerblank[2.4cm]{broadens}, so the
        energy is spread over more states.
\end{itemize}

\notesectionz{The four situations you must recognize}{17.1}

\begin{center}
\begin{tabular}{@{}lll@{}}
\hline
\textbf{Change} & \textbf{$\Delta S$} & \textbf{Why} \\
\hline
solid $\to$ liquid $\to$ gas & \textbf{positive} & particles freer, larger
  volume \\
gas expands (constant $T$) & \textbf{positive} & same particles, more space
  \\
temperature rises & \textbf{positive} & energy spread over more states \\
moles of gas increase & \textbf{positive} & the dominant term in a
  reaction \\
\hline
\end{tabular}
\end{center}

\begin{apctrap}
\textbf{For a reaction, count moles of gas first.} Gases carry far more
entropy than liquids or solids, so $\Delta n_{\text{gas}}$ almost always
decides the sign. Compare the standard entropies:
\ce{H2O(l)} is \SI{70}{\joule\per\kelvin\per\mole} but \ce{H2O(g)} is
\SI{189}{\joule\per\kelvin\per\mole} --- the same substance, nearly triple.

Only when $\Delta n_{\text{gas}} = 0$ do you look at anything else.
\end{apctrap}

\segment{GUIDED PRACTICE}{15}{Predicting the sign}

Give the sign of $\Delta S^\circ$ and the reason:

\begin{enumerate}[leftmargin=*,itemsep=0.35em]
  \item \ce{CaCO3(s) -> CaO(s) + CO2(g)}
        \answerblank[6cm]{positive --- 0 to 1 mol gas}
  \item \ce{N2(g) + 3H2(g) -> 2NH3(g)}
        \answerblank[6cm]{negative --- 4 to 2 mol gas}
  \item \ce{2H2(g) + O2(g) -> 2H2O(l)}
        \answerblank[6cm]{negative --- 3 mol gas to 0}
  \item \ce{H2O(s) -> H2O(l)}
        \answerblank[6cm]{positive --- solid to liquid}
  \item \ce{2NO2(g) -> N2O4(g)}
        \answerblank[6cm]{negative --- 2 to 1 mol gas}
\end{enumerate}

\segment{INSTRUCTION B}{20}{The second law}

\notesectionz{The universe, not the system}{17.2}
\practice{6}

\begin{center}
\fbox{\parbox{0.86\linewidth}{\centering
\textbf{Second law of thermodynamics:} in any thermodynamically favored
process, the entropy of the \emph{universe} increases. \\[0.3em]
$\Delta S_{\text{univ}} = \Delta S_{\text{sys}} +
\Delta S_{\text{surr}} > 0$}}
\end{center}

The system's entropy is allowed to \answerblank[2cm]{decrease} --- water
freezes, ammonia forms, life builds ordered structures --- provided the
surroundings gain more than the system loses.

The surroundings gain entropy when the system releases heat to them, and how
much they gain depends on temperature:
\[ \Delta S_{\text{surr}} = -\frac{\Delta H}{T} \]

\begin{apcnote}[Why the same heat matters more when it is cold]
The $T$ in the denominator is doing real work. Dumping
\SI{1}{\kilo\joule} into cold surroundings is a large fractional increase in
their energy dispersal; the same kilojoule into hot surroundings barely
registers.

That is why \textbf{exothermicity is a strong driving force at low
temperature and a weak one at high temperature} --- and it is the reason
water freezes below \SI{0}{\celsius} but melts above it, with no change in
$\Delta H$ or $\Delta S$ at all.
\end{apcnote}

\segment{APPLICATION}{20}{Reasoning with the second law}

\begin{enumerate}[leftmargin=*,itemsep=0.4em]
  \item Water freezing has $\Delta S_{\text{sys}} < 0$. Explain how it can
        still occur below \SI{0}{\celsius}.
        \answerlines[3]{Freezing is exothermic, so
        $\Delta S_{\text{surr}} = -\Delta H/T$ is positive. At low enough
        $T$ that positive term is larger in magnitude than the system's
        loss, so $\Delta S_{\text{univ}} > 0$ and the process is favored.
        Above \SI{0}{\celsius} the surroundings term is too small and the
        reverse process wins instead.}
  \item A living organism builds highly ordered molecules from simple ones.
        Does this violate the second law? \answerlines[3]{No. The organism
        is not an isolated system. It releases heat and simpler waste
        products to its surroundings, and the entropy those gain exceeds
        the entropy lost in building the ordered structures. The
        \emph{universe} still gains.}
\end{enumerate}

\begin{exitticket}
State the sign of $\Delta S_{\text{sys}}$, of $\Delta S_{\text{surr}}$, and
of $\Delta S_{\text{univ}}$ for the combustion of methane at room
temperature.
\keyonly{\begin{apcanswer}$\Delta S_{\text{sys}}$ is negative (3 mol gas
becomes 1 mol gas plus liquid water). $\Delta S_{\text{surr}}$ is strongly
positive, because the reaction is very exothermic. $\Delta S_{\text{univ}}$
is positive --- the reaction is thermodynamically favored, so the
surroundings term must dominate.\end{apcanswer}}
\end{exitticket}

% =====================================================================
\blocklesson{Calculating Entropy Changes}{Zumdahl \S17.6}

\begin{objectives}
  \item Calculate $\Delta S^\circ$ from absolute (standard molar) entropies.
  \item Explain why absolute entropies exist but absolute enthalpies do not.
\end{objectives}

\reading{Zumdahl \S17.6}{851--855}

\begin{warmup}[3]
  \item Sign of $\Delta S^\circ$ when moles of gas increase:
        \answerblank[2cm]{positive}
  \item $\Delta S_{\text{surr}} = $ \answerblank[2.6cm]{$-\Delta H/T$}
  \item Units of $S^\circ$: \answerblank[3.4cm]{J/(K$\cdot$mol)}
\end{warmup}

\segment{INSTRUCTION A}{25}{Absolute entropies}

\notesectionz{The third law makes a real zero possible}{17.6}
\practice{5}

The \term{third law of thermodynamics} says a perfect crystal at
\SI{0}{\kelvin} has an entropy of exactly \answerblank[1.4cm]{zero}. That
gives entropy a genuine origin, so every substance has an
\term{absolute entropy} $S^\circ$ --- a positive number, tabulated
directly.

\begin{apctrap}
Contrast this with enthalpy. There is no absolute $H$, so tables list
$\Delta H_f^\circ$, the enthalpy \emph{of formation}, and elements in their
standard states are assigned zero by convention.

Entropy tables are different: \textbf{$S^\circ$ for an element is not zero}.
\ce{O2(g)} has $S^\circ = \SI{205}{\joule\per\kelvin\per\mole}$. Setting
element entropies to zero out of habit is a reliable way to lose a point.
\end{apctrap}

\[ \boxed{\;\Delta S^\circ_{\text{reaction}} = \sum n S^\circ_{\text{products}}
   - \sum n S^\circ_{\text{reactants}}\;} \]

Note the coefficients multiply, exactly as in Hess's law --- and note the
units: $S^\circ$ is in \textbf{joules} while $\Delta H^\circ$ is in
\textbf{kilojoules}. Mixing them is the single most common arithmetic error
in this unit.

\begin{workedexample}[Worked example 1: decomposing limestone]
\[ \ce{CaCO3(s) -> CaO(s) + CO2(g)} \]
$S^\circ$: \ce{CaCO3} 93, \ce{CaO} 40, \ce{CO2} 214
\si{\joule\per\kelvin\per\mole}.

\[ \Delta S^\circ = (40 + 214) - (93) = \mathbf{+161}~
   \si{\joule\per\kelvin} \]

Positive, as predicted --- a gas appeared where there was none.
\end{workedexample}

\begin{workedexample}[Worked example 2: the Haber process]
\[ \ce{N2(g) + 3H2(g) -> 2NH3(g)} \]
$S^\circ$: \ce{N2} 192, \ce{H2} 131, \ce{NH3} 193.

\[ \Delta S^\circ = 2(193) - [192 + 3(131)]
   = 386 - 585 = \mathbf{-199}~\si{\joule\per\kelvin} \]

Four moles of gas became two, so the sign had to be negative. Notice the
coefficient of 3 on hydrogen --- dropping it gives $-68$ and the wrong
magnitude entirely.
\end{workedexample}

\segment{GUIDED PRACTICE}{15}{Running the sums}

$S^\circ$ values (\si{\joule\per\kelvin\per\mole}): \ce{H2} 131,
\ce{O2} 205, \ce{H2O(l)} 70, \ce{H2O(g)} 189, \ce{NO2} 240,
\ce{N2O4} 304.

\begin{enumerate}[leftmargin=*,itemsep=0.4em]
  \item \ce{2H2(g) + O2(g) -> 2H2O(l)}: \workspace[2cm]
        \keyonly{\begin{apcanswer}$2(70) - [2(131) + 205] = 140 - 467
        = \mathbf{-327}~\si{\joule\per\kelvin}$\end{apcanswer}}
  \item \ce{H2O(l) -> H2O(g)}: \answerblank[3.4cm]{$+119$ J/K}
  \item \ce{2NO2(g) -> N2O4(g)}: \answerblank[3.4cm]{$-176$ J/K}
\end{enumerate}

\segment{APPLICATION}{20}{Interpreting the numbers}

\begin{enumerate}[leftmargin=*,itemsep=0.4em]
  \item For \ce{H2O(l) -> H2O(g)}, $\Delta S^\circ = +119$~J/K while
        $\Delta H^\circ = +44$~kJ. Explain what each sign says about
        whether the process is favored at \SI{25}{\celsius}.
        \answerlines[3]{The entropy term favors vaporization --- gas is far
        more dispersed. The enthalpy term opposes it, since hydrogen bonds
        must be broken. At \SI{25}{\celsius} the enthalpy term wins and
        liquid water does not spontaneously all become vapour; at higher
        temperature the $T\Delta S$ term grows until it does.}
  \item Why is $S^\circ$ for \ce{H2O(g)} nearly triple that for
        \ce{H2O(l)}, when it is the same substance?
        \answerlines[2]{In the gas the molecules are free to occupy a far
        larger volume and to move independently, so vastly more microscopic
        arrangements correspond to the same macroscopic state.}
  \item A student computes $\Delta S^\circ$ for the Haber process as
        $193 - (192 + 131) = -130$~J/K. Find both errors.
        \answerlines[3]{They omitted the coefficient 2 on \ce{NH3} and the
        coefficient 3 on \ce{H2}. Correctly,
        $2(193) - [192 + 3(131)] = -199$~J/K.}
\end{enumerate}

\begin{exitticket}
Why can a table list an absolute $S^\circ$ for oxygen gas but not an
absolute $H^\circ$?
\keyonly{\begin{apcanswer}The third law fixes a true zero for entropy --- a
perfect crystal at \SI{0}{\kelvin} has $S = 0$ --- so entropies can be
measured on an absolute scale. Enthalpy has no such natural zero, so only
\emph{changes} in enthalpy are measurable and tables use
$\Delta H_f^\circ$ with elements assigned zero by
convention.\end{apcanswer}}
\end{exitticket}

% =====================================================================
\blocklesson{Gibbs Free Energy}{Zumdahl \S17.3--17.4, 17.7}

\begin{objectives}
  \item Calculate $\Delta G^\circ$ two ways and judge thermodynamic
        favorability.
  \item Predict the temperature conditions under which a process is favored.
\end{objectives}

\reading{Zumdahl \S17.3--17.4, 17.7}{843--851, 855--860}

\begin{warmup}[4]
  \item $\Delta S^\circ$ for \ce{CaCO3(s) -> CaO(s) + CO2(g)}:
        \answerblank[2.6cm]{$+161$ J/K}
  \item $\Delta S_{\text{univ}}$ must be what sign for a favored process?
        \answerblank[2cm]{positive}
  \item Units of $S^\circ$ versus $\Delta H^\circ$:
        \answerblank[5.2cm]{J/(K$\cdot$mol) versus kJ/mol}
  \item The AP term replacing ``spontaneous'':
        \answerblank[5.4cm]{thermodynamically favored}
\end{warmup}

\segment{INSTRUCTION A}{25}{One number that decides}

\notesectionz{Folding the surroundings into the system}{17.4}
\practice{5}

Judging favorability from $\Delta S_{\text{univ}}$ means tracking the
surroundings, which is a nuisance. \term{Gibbs free energy} removes the
need:
\[ G = H - TS \qquad\Longrightarrow\qquad
   \boxed{\;\Delta G^\circ = \Delta H^\circ - T\Delta S^\circ\;} \]

Dividing $\Delta G = \Delta H - T\Delta S$ by $-T$ gives exactly
$\Delta S_{\text{surr}} + \Delta S_{\text{sys}} = \Delta S_{\text{univ}}$.
So $\Delta G$ is the second law rewritten \emph{entirely in terms of the
system} --- which is why it is worth having.

\begin{center}
\fbox{\parbox{0.8\linewidth}{\centering
$\Delta G^\circ < 0$: \textbf{thermodynamically favored} \\
$\Delta G^\circ > 0$: not favored (the reverse is) \\
$\Delta G^\circ = 0$: at equilibrium}}
\end{center}

The sign flip is worth stating: $\Delta S_{\text{univ}}$ must be
\answerblank[1.8cm]{positive}, but $\Delta G^\circ$ must be
\answerblank[1.8cm]{negative}, because of the $-T$ used in the derivation.

\notesectionz{The second route: free energies of formation}{17.7}

Exactly as with enthalpy, $\Delta G^\circ$ can be assembled from tabulated
formation values:
\[ \Delta G^\circ_{\text{reaction}} =
   \sum n\Delta G_f^\circ{}_{\text{products}}
   - \sum n\Delta G_f^\circ{}_{\text{reactants}} \]

\begin{workedexample}[Worked example 3: two routes, one answer]
For \ce{2H2(g) + O2(g) -> 2H2O(l)}:

\textbf{Route 1 --- from $\Delta H^\circ$ and $\Delta S^\circ$.}
$\Delta H^\circ = 2(-286) = -572$~kJ;
$\Delta S^\circ = -327$~J/K $= -0.327$~kJ/K.
\[ \Delta G^\circ = -572 - (298)(-0.327) = -572 + 97 = \mathbf{-475}~
   \text{kJ} \]

\textbf{Route 2 --- from $\Delta G_f^\circ$.}
$\Delta G^\circ = 2(-237) - 0 = \mathbf{-474}$~kJ.

The two agree to within rounding in the table. \textbf{Watch the units} ---
the entropy had to be converted from J/K to kJ/K before it could be
subtracted from a value in kJ.
\end{workedexample}

\segment{INSTRUCTION B}{20}{When temperature decides}

\notesectionz{The four cases}{17.3}
\practice{5}

Because $\Delta G^\circ = \Delta H^\circ - T\Delta S^\circ$, the signs of
$\Delta H^\circ$ and $\Delta S^\circ$ between them determine whether
temperature matters at all:

\begin{center}
\begin{tabular}{@{}cccl@{}}
\hline
$\Delta H^\circ$ & $\Delta S^\circ$ & \textbf{Favored at} &
  \textbf{Comment} \\
\hline
$-$ & $+$ & \textbf{all temperatures} & both terms help; no calculation
  needed \\
$+$ & $-$ & \textbf{no temperature} & both terms oppose; never favored \\
$+$ & $+$ & \textbf{high $T$} & entropy wins once $T$ is large enough \\
$-$ & $-$ & \textbf{low $T$} & enthalpy wins while $T$ is small \\
\hline
\end{tabular}
\end{center}

For the two mixed cases, the \term{crossover temperature} --- where
$\Delta G^\circ = 0$ and the behaviour switches --- is
\[ T = \frac{\Delta H^\circ}{\Delta S^\circ} \]

\begin{workedexample}[Worked example 4: why limestone needs a kiln]
\ce{CaCO3(s) -> CaO(s) + CO2(g)}:
$\Delta H^\circ = +178.5$~kJ, $\Delta S^\circ = +161$~J/K.

Both positive $\Rightarrow$ favored only at high temperature.

\textbf{At \SI{298}{\kelvin}:}
$\Delta G^\circ = 178.5 - 298(0.161) = +131$~kJ --- \emph{not} favored, which
is why limestone buildings do not crumble into quicklime.

\textbf{Crossover:}
\[ T = \frac{178.5}{0.161} = \mathbf{1109}~\si{\kelvin}
   \quad (\approx \SI{836}{\celsius}) \]
Above about \SI{1100}{\kelvin} the decomposition becomes favored --- which is
exactly the temperature at which lime kilns are operated.
\end{workedexample}

\segment{APPLICATION}{20}{Predicting and calculating}

\begin{enumerate}[leftmargin=*,itemsep=0.4em]
  \item For the Haber process, $\Delta H^\circ = -92$~kJ and
        $\Delta S^\circ = -199$~J/K. Find $\Delta G^\circ$ at
        \SI{298}{\kelvin} and the crossover temperature. \workspace[2.8cm]
        \keyonly{\begin{apcanswer}$\Delta G^\circ = -92 - 298(-0.199)
        = -92 + 59 = \mathbf{-33}$~kJ, so favored at \SI{298}{\kelvin}.
        Crossover $T = 92/0.199 = \mathbf{462}$~K --- favored
        \emph{below} 462 K, since both terms are
        negative.\end{apcanswer}}
  \item For \ce{H2O(l) -> H2O(g)}, $\Delta H^\circ = +44$~kJ and
        $\Delta S^\circ = +119$~J/K. Find the crossover temperature and
        say what it physically is. \workspace[2.4cm]
        \keyonly{\begin{apcanswer}$T = 44/0.119 = \mathbf{370}$~K
        $\approx \SI{97}{\celsius}$ --- the \textbf{boiling point}. Above
        it, vaporization is favored at 1 atm; below it, condensation is.
        The tabulated value is 373 K; the 3 K gap is rounding in the
        table.\end{apcanswer}}
  \item Without calculating, say whether
        \ce{2H2O2(l) -> 2H2O(l) + O2(g)} ($\Delta H^\circ < 0$) is favored,
        and at what temperatures. \answerlines[2]{$\Delta H^\circ$ is
        negative and $\Delta S^\circ$ is positive (a gas is produced from a
        liquid), so both terms help and it is favored at \emph{all}
        temperatures. No calculation is needed.}
\end{enumerate}

\begin{exitticket}
A reaction has $\Delta H^\circ = +50$~kJ and $\Delta S^\circ = +100$~J/K.
Is it favored at \SI{298}{\kelvin}? Above what temperature does that change?
\keyonly{\begin{apcanswer}At \SI{298}{\kelvin},
$\Delta G^\circ = 50 - 298(0.100) = +20$~kJ, so it is not favored.
Crossover at $T = 50/0.100 = \SI{500}{\kelvin}$; above \SI{500}{\kelvin} it
becomes favored.\end{apcanswer}}
\end{exitticket}

% =====================================================================
\blocklesson{Free Energy, Equilibrium, and Kinetic Control}{Zumdahl \S17.9}

\begin{objectives}
  \item Relate $\Delta G^\circ$ to $K$ and estimate one from the other.
  \item Explain why a favored reaction may not occur at a measurable rate.
\end{objectives}

\reading{Zumdahl \S17.9}{863--869}

\begin{warmup}[4]
  \item $\Delta G^\circ = $ \answerblank[4cm]{$\Delta H^\circ -
        T\Delta S^\circ$}
  \item Favored means $\Delta G^\circ$ is: \answerblank[2cm]{negative}
  \item If $K > 1$, which side is favored?
        \answerblank[2.4cm]{products}
  \item What raises a rate without changing $K$?
        \answerblank[2.6cm]{a catalyst}
\end{warmup}

\segment{INSTRUCTION A}{25}{Two languages for the same fact}

\notesectionz{$\Delta G^\circ$ and $K$}{17.9}
\practice{5}

You now have two ways of saying ``products are favored'': $K > 1$ from
Unit 7, and $\Delta G^\circ < 0$ from this unit. They are the same
statement, linked by
\[ \boxed{\;\Delta G^\circ = -RT\ln K\;}
   \qquad\text{equivalently}\qquad K = e^{-\Delta G^\circ/RT} \]
with $R = \SI{8.314}{\joule\per\mole\per\kelvin}$ --- note \textbf{joules},
so $\Delta G^\circ$ must be converted from kJ.

\begin{center}
\begin{tabular}{@{}ccl@{}}
\hline
$\Delta G^\circ$ & $K$ & \\
\hline
negative & $> 1$ & products favored at equilibrium \\
zero     & $= 1$ & neither side favored \\
positive & $< 1$ & reactants favored \\
\hline
\end{tabular}
\end{center}

\begin{apcnote}[Estimating without a calculator]
The CED asks you to connect $K$ and $\Delta G^\circ$ \emph{qualitatively by
estimation}, so learn the scale rather than only the formula. The quantity
to compare against is $RT$, which at \SI{298}{\kelvin} is about
\SI{2.5}{\kilo\joule\per\mole}.

\begin{center}
\begin{tabular}{@{}rl@{}}
\hline
$\Delta G^\circ \approx 0$ & $K$ close to 1 \\
$\Delta G^\circ = -10$ kJ & $K \approx 60$ \\
$\Delta G^\circ = -50$ kJ & $K \approx 6\times10^{8}$ \\
$\Delta G^\circ = +50$ kJ & $K \approx 2\times10^{-9}$ \\
\hline
\end{tabular}
\end{center}

The lesson: because $\Delta G^\circ$ sits inside an exponential, a modest
free-energy change produces an enormous swing in $K$.
\end{apcnote}

\begin{workedexample}[Worked example 5: limestone again, as an equilibrium]
At \SI{298}{\kelvin} the decomposition of \ce{CaCO3} has
$\Delta G^\circ = +131$~kJ. Find $K$.

\[ \ln K = -\frac{\Delta G^\circ}{RT}
   = -\frac{131{,}000}{(8.314)(298)} = -52.9 \]
\[ K = e^{-52.9} = \mathbf{1\times10^{-23}} \]

Essentially no \ce{CO2} above limestone at room temperature --- which is
what ``not favored'' looks like quantitatively.
\end{workedexample}

\segment{INSTRUCTION B}{20}{Favored is not the same as fast}

\notesectionz{Kinetic control}{17.4 (CED)}
\practice{6}

This is CED topic \textbf{9.4}, and it is the payoff for having taught
kinetics first.

A thermodynamically favored reaction tells you \emph{where} the system would
end up. It says nothing about \emph{how long} that takes. When a favored
reaction proceeds too slowly to observe, it is under
\term{kinetic control}, and the usual cause is a high
\answerblank[3.8cm]{activation energy}.

\begin{apctrap}
\textbf{A reaction that is not proceeding is not necessarily at
equilibrium.} Diamond converting to graphite has $\Delta G^\circ < 0$ at
room temperature --- it is favored --- yet diamonds persist indefinitely
because the activation barrier is enormous.

If you are told a process is favored but nothing appears to happen, the
conclusion the CED wants is \textbf{kinetic control}, not equilibrium.

Note what a catalyst does and does not do: it lowers $E_a$ and so raises the
rate, but it cannot change $\Delta G^\circ$ or $K$. Catalysts move a system
to equilibrium faster; they never move the equilibrium.
\end{apctrap}

\segment{APPLICATION}{20}{Thermodynamics versus kinetics}

\begin{enumerate}[leftmargin=*,itemsep=0.4em]
  \item A reaction has $\Delta G^\circ = -33$~kJ at \SI{298}{\kelvin}.
        Calculate $K$. \workspace[2.2cm]
        \keyonly{\begin{apcanswer}$\ln K = 33{,}000/[(8.314)(298)] = 13.3$;
        $K = e^{13.3} = \mathbf{6\times10^{5}}$ --- strongly
        product-favored\end{apcanswer}}
  \item A mixture of \ce{H2} and \ce{O2} can sit in a flask for years
        without reacting, yet the combustion has
        $\Delta G^\circ = -474$~kJ. Explain, and say what a spark does.
        \answerlines[3]{The reaction is strongly favored but under kinetic
        control: the activation energy is very high, so at room temperature
        almost no collisions have enough energy. A spark supplies that
        energy locally and starts the reaction, which then releases enough
        heat to sustain itself. The spark changes the rate, not
        $\Delta G^\circ$.}
  \item A student says ``$\Delta G^\circ$ is negative, so the reaction will
        be fast.'' Correct them in one sentence.
        \answerlines[2]{$\Delta G^\circ$ predicts the position of
        equilibrium, not the rate --- rate is set by activation energy, and
        a strongly favored reaction can be immeasurably slow.}
\end{enumerate}

\begin{exitticket}
Give one reaction that is thermodynamically favored but does not visibly
occur, and name what is preventing it.
\keyonly{\begin{apcanswer}Diamond $\to$ graphite (or the combustion of a
hydrogen--oxygen mixture without ignition). In each case
$\Delta G^\circ < 0$, but a very high activation energy puts the process
under kinetic control, so the rate is negligible at room
temperature.\end{apcanswer}}
\end{exitticket}

% =====================================================================
\blocklesson{Dissolution and Coupled Reactions}{Zumdahl \S17.5, 17.10}

\begin{objectives}
  \item Explain solubility in terms of the enthalpy and entropy of
        dissolution.
  \item Explain how an unfavorable process can be driven by coupling or an
        external energy source.
\end{objectives}

\reading{Zumdahl \S17.5, 17.10}{850--851, 869--873}

\begin{warmup}[3]
  \item $\Delta G^\circ = -RT\ln K$; if $\Delta G^\circ < 0$ then $K$ is:
        \answerblank[3cm]{greater than 1}
  \item A favored but immeasurably slow reaction is under:
        \answerblank[3.4cm]{kinetic control}
  \item Adding two reactions adds their:
        \answerblank[3.4cm]{$\Delta G^\circ$ values}
\end{warmup}

\segment{INSTRUCTION A}{25}{Why some salts dissolve and others do not}

\notesectionz{Three competing contributions}{17.5}
\practice{4}

This is CED \textbf{9.6}. Dissolving a salt is a free-energy balance with
three parts:

\begin{enumerate}[leftmargin=*,itemsep=0.3em]
  \item \textbf{Breaking up the solid} --- overcoming the lattice.
        $\Delta H$ strongly \answerblank[2cm]{positive},
        $\Delta S$ positive.
  \item \textbf{Reorganizing the solvent} around the ions.
        $\Delta S$ \answerblank[2cm]{negative} --- water molecules become
        ordered in hydration shells.
  \item \textbf{Ion--solvent attraction} (hydration).
        $\Delta H$ \answerblank[2cm]{negative}.
\end{enumerate}

\begin{apcnote}[Why the CED does not ask you to predict the total]
The framework says outright that predicting the overall $\Delta G^\circ$ of
dissolution ``can be challenging due to the cancellations among'' these
three factors. The enthalpy terms are large and opposite; the entropy terms
likewise. What is left over is a small difference between big numbers.

So you are expected to \textbf{identify the contributions and their signs},
not to predict solubility from scratch. An answer that reasons through the
three factors earns credit even if the final direction is uncertain.
\end{apcnote}

\begin{workedexample}[Worked example 6: dissolving that cools the beaker]
\ce{NH4NO3} dissolves \emph{endothermically} --- the beaker gets cold, so
$\Delta H^\circ > 0$. Yet it is very soluble. Why?

Because $\Delta S^\circ$ is strongly positive: an ordered crystal becomes
freely moving hydrated ions dispersed through the solvent. At room
temperature $T\Delta S^\circ$ exceeds $\Delta H^\circ$, so
$\Delta G^\circ = \Delta H^\circ - T\Delta S^\circ < 0$.

\textbf{Entropy is driving this one on its own} --- the enthalpy term is
working against it. This is precisely the case the CED flags in 9.3.A.4,
where you must consider both terms rather than assuming exothermic means
favored.
\end{workedexample}

\segment{INSTRUCTION B}{20}{Driving the unfavorable}

\notesectionz{Coupling, and external energy}{17.10}
\practice{4}

A process with $\Delta G^\circ > 0$ will not happen on its own. There are
two ways to make it happen anyway --- CED topic \textbf{9.7}.

\textbf{1. Supply external energy.} Examples the CED names:

\begin{itemize}[leftmargin=*,itemsep=0.25em]
  \item \textbf{Electrical energy} driving an electrolytic cell or charging
        a battery --- the whole subject of Chapter 18.
  \item \textbf{Light} driving photosynthesis, which converts \ce{CO2} and
        water into glucose with a hugely positive $\Delta G^\circ$.
\end{itemize}

\textbf{2. Couple it to a favorable reaction.} Because free energies
\answerblank[2cm]{add}, an unfavorable step can be carried by a favorable
one, provided the two share a \answerblank[4.4cm]{common intermediate} and
the \emph{sum} has $\Delta G^\circ < 0$.

\begin{workedexample}[Worked example 7: how cells pay for chemistry]
Phosphorylating glucose is unfavorable on its own:
\[ \text{glucose} + \ce{P_i} \to \text{glucose-6-phosphate}
   \qquad \Delta G^\circ = +13.8~\text{kJ/mol} \]
Hydrolysing ATP is strongly favorable:
\[ \ce{ATP} + \ce{H2O} \to \ce{ADP} + \ce{P_i}
   \qquad \Delta G^\circ = -30.5~\text{kJ/mol} \]
Coupled through the shared \ce{P_i}, the overall process is
\[ \Delta G^\circ = +13.8 + (-30.5) = \mathbf{-16.7}~\text{kJ/mol} \]
and therefore favored. This is what ``ATP is the cell's energy currency''
actually means --- not that ATP contains energy, but that its hydrolysis is
favorable enough to pay for reactions that are not.
\end{workedexample}

\segment{APPLICATION}{20}{Putting the unit together}

\begin{enumerate}[leftmargin=*,itemsep=0.4em]
  \item Photosynthesis has $\Delta G^\circ \approx +2870$~kJ/mol of
        glucose. Explain how it occurs at all.
        \answerlines[3]{It is driven by an external energy source ---
        sunlight. The absorbed light supplies the free energy the reaction
        lacks, so the overall system including the light source still has
        $\Delta G < 0$. Nothing about the second law is violated.}
  \item Reaction X has $\Delta G^\circ = +25$~kJ. It is coupled to
        reaction Y with $\Delta G^\circ = -60$~kJ. Is the coupled process
        favored, and what must the two reactions share?
        \answerblank[7.4cm]{yes, $-35$ kJ; a common intermediate}
  \item \ce{NaCl} dissolves with $\Delta H^\circ$ very close to zero.
        Explain what must be driving it. \answerlines[3]{The entropy term.
        With $\Delta H^\circ \approx 0$, $\Delta G^\circ \approx
        -T\Delta S^\circ$, and dissolution disperses an ordered lattice
        into mobile hydrated ions, so $\Delta S^\circ > 0$ and
        $\Delta G^\circ < 0$.}
\end{enumerate}

\begin{exitticket}
A classmate says ``an unfavorable reaction can never happen.'' Give the two
ways it can, with one example each.
\keyonly{\begin{apcanswer}It can be driven by an external energy source ---
electricity in electrolysis, or light in photosynthesis. Or it can be
coupled to a favorable reaction sharing a common intermediate, so that the
summed $\Delta G^\circ$ is negative --- as ATP hydrolysis drives the
phosphorylation of glucose.\end{apcanswer}}
\end{exitticket}

\end{document}
